\documentclass[12pt]{article}

\usepackage{amsmath, amsfonts, amssymb, enumerate}
\usepackage[margin=1cm]{geometry}

\begin{document}

\noindent
{\large\bfseries Number Theory}
\hfill
znz
\quad
\today

\hrule
\medskip

\begin{enumerate}
  \item
        Let
        \[
        f,g\in F[t],
        \qquad
        g\neq0.
        \]Prove that there exist unique
        \[
        q,r\in F[t]
        \]such that
        \[
        f=qg+r
        \]and either
        \[
        r=0
        \]or
        \[
        \deg r<\deg g.
        \]Your proof must establish both:
        \[
        \text{existence}
        \]and
        \[
        \text{uniqueness}.
        \]
        Do not assume polynomial long division as an algorithmic fact; derive the theorem from the field structure and degree.
  \item
        Let
        \[
        f,g\in F[t]
        \]be not both zero.
        Prove that there exists a unique monic polynomial \(d\) such that
        \[
        d\mid f,\qquad d\mid g,
        \]and
        \[
        c\mid f,\ c\mid g
        \implies
        c\mid d.
        \]Then prove that there exist
        \[
        x,y\in F[t]
        \]such that
        \[
        \boxed{d=xf+yg.}
        \]Your argument should ultimately depend on Problem 1.
        As part of your proof, justify why the Euclidean algorithm must terminate.
  \item
        Let
        \[
        \pi\in F[t]
        \]be irreducible.

        \begin{enumerate}[(a)]
          \item
                Prove that if
                \[
                \pi\nmid f,
                \]then
                \[
                \gcd(\pi,f)=1.
                \]
          \item
                Using Problem 2, prove that
                \[
                \pi\mid fg
                \]implies
                \[
                \pi\mid f
                \]or
                \[
                \pi\mid g.
                \]Thus determine the relation between irreducibility and primality in \(F[t]\).
          \item
                Prove that every nonzero nonunit polynomial in \(F[t]\) can be written as a finite product of irreducible polynomials.

          \item
                Prove uniqueness:
                if
                \[
                f
                =
                u\pi_1\cdots\pi_r
                =
                v\rho_1\cdots\rho_s,
                \]where
                \[
                u,v\in F^\times
                \]and all \(\pi_i,\rho_j\) are irreducible, prove that
                \[
                r=s
                \]and, after reordering, each \(\pi_i\) is associate to \(\rho_i\).
        \end{enumerate}
        Conclude that every monic nonconstant polynomial has a unique factorization into monic irreducible polynomials, up to ordering.
  \item
        Let
        \[
        m\in F[t],
        \qquad
        \deg m\ge1.
        \]Prove the equivalence
        \[
        \boxed{
        F[t]/(m)\text{ is a field}
        \iff
        m\text{ is irreducible}.
        }
        \]You must prove both directions.
        For the irreducible direction, begin with an arbitrary nonzero class
        \[
        [f]\neq[0].
        \]For the reducible direction, reason from a nontrivial factorization
        \[
        m=ab.
        \]Your proof should use results established earlier in today's sheet rather than treating this as an independent theorem.
  \item
        Let
        \[
        p
        \]be prime and let
        \[
        m\in\mathbb F_p[t]
        \]be irreducible of degree
        \[
        d\ge1.
        \]Set
        \[
        K=\mathbb F_p[t]/(m).
        \]

        \begin{enumerate}[(a)]
          \item
                Prove that every element of \(K\) has a unique representative
                \[
                a_0+a_1t+\cdots+a_{d-1}t^{d-1}
                \]with
                \[
                a_i\in\mathbb F_p.
                \]
          \item
                Deduce that
                \[
                \boxed{|K|=p^d.}
                \]
          \item
                Let
                \[
                \alpha=[t]\in K.
                \]If
                \[
                m(t)
                =
                t^d+c_{d-1}t^{d-1}+\cdots+c_0,
                \]prove that
                \[
                m(\alpha)=0
                \]inside \(K\).
                In other words,
                \[
                \alpha^d
                +
                c_{d-1}\alpha^{d-1}
                +\cdots+c_0
                =
                0.
                \]
          \item
                Prove that every element of \(K\) can be written uniquely as
                \[
                a_0+a_1\alpha+\cdots+a_{d-1}\alpha^{d-1}.
                \]
          \item
                Explain precisely why \(K\) contains a copy of \(\mathbb F_p\), but when \(d>1\),
                \[
                K\neq\mathbb F_p.
                \]At this point you will have constructed, from first principles, a field with
                \[
                p^d
                \]elements.
        \end{enumerate}

\end{enumerate}

\end{document}
